\documentclass[polish, 10pt]{extarticle}
\usepackage[top=2cm, bottom=2cm, left=3cm, right=3cm]{geometry}
\usepackage{graphicx} % Required for inserting images

\usepackage[polish]{babel}
\usepackage{polski}
\usepackage{placeins}
\usepackage{dirtytalk}
\usepackage{csquotes}
\usepackage{float}
\usepackage{amsmath}
\usepackage{graphicx}
\usepackage{hyperref}
\usepackage{graphicx}
\usepackage{listings}
\usepackage{xcolor}
\usepackage{amssymb}
\usepackage{enumitem}
\usepackage{float}
\definecolor{codegreen}{rgb}{0,0.6,0}
\definecolor{codegray}{rgb}{0.5,0.5,0.5}
\definecolor{codepurple}{rgb}{0.58,0,0.82}
\definecolor{backcolour}{rgb}{0.95,0.95,0.92}

\lstdefinestyle{mystyle}{
    backgroundcolor=\color{backcolour},   
    commentstyle=\color{codegreen},
    keywordstyle=\color{magenta},
    numberstyle=\tiny\color{codegray},
    stringstyle=\color{codepurple},
    basicstyle=\ttfamily\footnotesize,
    breakatwhitespace=false,         
    breaklines=true,                 
    captionpos=b,                    
    keepspaces=true,                 
    numbers=left,                    
    numbersep=5pt,                  
    showspaces=false,                
    showstringspaces=false,
    showtabs=false,                  
    tabsize=2
}
\lstset{style=mystyle}
\hypersetup{
    colorlinks=true,
    linktoc=all     %set to all if you want both sections and subsections linked
}
\usepackage{mdframed}
\usepackage{titling}
\setlength{\droptitle}{-2cm} 

\begin{document}

\title{Rachunek Prawdopodobieństwa i Statystyka -- Kolokwium II}
\author{}
\date{}

\maketitle
\vspace{-2cm}

\section*{Zadanie 1 (30 pkt)}
Niech zmienna losowa dwuwymiarowa \((X, Y)\) ma rozkład ciągły o gęstości
\[
f(x, y) =
\begin{cases}
    c x^2 y & \text{jeśli } x < 0, \, y < 1, \, x + y > 0, \\
    0 & \text{w przeciwnym wypadku},
\end{cases}
\]
dla pewnej stałej \(c \in \mathbb{R}\). Wyznacz wartość \(c\). Wyznacz obie gęstości brzegowe i obie gęstości warunkowe. Oblicz prawdopodobieństwo
\[
P\left(X > -\frac{1}{4} \, \middle| \, Y = \frac{1}{3} \right),
\]
zbadaj czy zmienne $X$ i $Y$ są niezależne.

\section*{Zadanie 2 (15 pkt)}
% TODO: zmienić wariancję na średnie
\noindent Poniżej znajduje się raport z testu statystycznego. Zakładamy, że dane pochodzą z dwóch rozkładów normalnych.

\begin{enumerate}
    \item[a)] Podaj definicję statystyki testowej używanej w tym teście.
    \item[b)] Opisz rozkład prawdopodobieństwa tej statystyki. Jakie są długości obu badanych prób?
    \item[c)] Zinterpretuj wyniki testu.
\end{enumerate}

\begin{verbatim}
## F test to compare two variances
## data: x and y
## F = 0.6123, num df = 119, denom df = 99, p-value = 0.01058
## alternative hypothesis: true ratio of variances is not equal to 1
## 95 percent confidence interval:
##  0.4177287 0.8915407
## sample estimates:
## ratio of variances
## 0.612303
\end{verbatim}

\section*{Zadanie 3 (25 pkt)}
\begin{enumerate}
    \item[a)] Wyznacz estymator największej wiarygodności w modelu statystycznym $\{ \text{Geo}(p) : p \in (0, 1) \}$.

    \item[b)] Wyznacz estymator parametru $\lambda$ rozkładu wykładniczego $\text{Exp}(\lambda)$ przy użyciu metody momentów, tzn. w modelu statystycznym: $\{ \text{Exp}(\lambda) : \lambda > 0 \}$.

\end{enumerate}
\section*{Zadanie 4 (15 pkt)}
Zmierzono czas życia, czyli czas działania, próby losowej 16 żarówek o ustalonej mocy. 
Średni czas życia w próbie wyniósł \( 3000 \) godzin, natomiast odchylenie standardowe \( 20 \) godzin. 
Przy założeniu, że czas życia żarówki jest zmienną losową o rozkładzie normalnym, podać przedział ufności dla wartości średniej tego rozkładu na poziomie ufności \( 0.9 \). Opisz sposób konstrukcji przedziału ufności w tym przypadku.

\section*{Zadanie 5 (15 pkt)}
Niech $X$ będzie zmienną losową. Wykaż, że jeśli $Y = aX + b$ dla pewnych $a, b \in \mathbb{R}$ (i $a \neq 0$) to $|\rho_{X,Y}| = 1$.
\end{document}
